\kbvsechead{ಒತ್ತಕ್ಷರಗಳು :}

\indent ಎರಡು ಅಥವಾ ಮೂರು ವ್ಯಂಜನಗಳನ್ನು ಒತ್ತಿ ಕೂಡಿಸಿ ಒಂದಕ್ಷರವಾಗಿಸಿ ಉಚ್ಚರಿಸಬಹುದು. ಈ ಬಗೆಯ ಸಂಯುಕ್ತಾಕ್ಷರಗಳಿಗೆ ಒತ್ತಕ್ಷರಗಳೆಂದು ಹೆಸರು.\par

\indent ಉದಾಹರಣೆಗೆ,\par

\begin{kbvderiv}
\vspace{4pt plus 3pt minus 1pt}
\begin{center}\begin{adjustbox}{max width=\textwidth}
{\renewcommand{\arraystretch}{1.22}\setlength{\tabcolsep}{4pt}
\begin{tabular}{l c l}
ಟ್ + ಟ್ + (ಅ) & {\color{kbvgrey}$=$} & \kbvresult{ಟ್ಟ} \\
ಸ್ + ತ್ + ರ್ + (ಅ) & {\color{kbvgrey}$=$} & \kbvresult{ಸ್ತ್ರ} \\
\end{tabular}}
\end{adjustbox}\end{center}
\vspace{4pt plus 3pt minus 1pt}
\end{kbvderiv}

\begin{kbvrule}
\noindent\kbvbullet\hspace{0.5em}ಒಂದೇ ಜಾತಿಯ ವ್ಯಂಜನಗಳು ಸೇರಿ ಆಗುವ ಒತ್ತಕ್ಷರಗಳಿಗೆ ಸಜಾತೀಯ ಒತ್ತಕ್ಷರಗಳೆಂದು ಹೆಸರು.
\end{kbvrule}

\vspace{4pt plus 3pt minus 1pt}
\begin{center}
{\renewcommand{\arraystretch}{1.28}\setlength{\tabcolsep}{11pt}
\begin{tabular}{c c c c c}
ಕ್ಕ & ಖ್ಖ & ಗ್ಗ & ಘ್ಘ & ಙ್ಙ \\
ಚ್ಚ & ಛ್ಛ & ಜ್ಜ & ಝ್ಝ & ಞ್ಞ \\
ಟ್ಟ & ಠ್ಠ & ಡ್ಡ & ಢ್ಢ & ಣ್ಣ \\
ತ್ತ & ಥ್ಥ & ದ್ದ & ಧ್ಧ & ನ್ನ \\
ಪ್ಪ & ಫ್ಫ & ಬ್ಬ & ಭ್ಭ & ಮ್ಮ \\
\end{tabular}}
\end{center}
\vspace{2pt}

\begin{center}{\fontsize{11pt}{15pt}\selectfont ಯ್ಯ\hspace{1.1em}ರ್ರ\hspace{1.1em}ಲ್ಲ\hspace{1.1em}ವ್ವ\hspace{1.1em}ಶ್ಶ\hspace{1.1em}ಷ್ಷ\hspace{1.1em}ಸ್ಸ\hspace{1.1em}ಹ್ಹ\hspace{1.1em}ಳ್ಳ}\end{center}\vspace{3pt}

\noindent\hspace*{2.2em}—ಇವು ಸಜಾತೀಯ ಒತ್ತಕ್ಷರಗಳು.\par

\begin{kbvrule}
\noindent\kbvbullet\hspace{0.5em}ವಿಭಿನ್ನ ಜಾತಿಯ ವ್ಯಂಜನಗಳು ಸೇರಿ ಆಗುವ ಒತ್ತಕ್ಷರಗಳಿಗೆ ವಿಜಾತೀಯ ಒತ್ತಕ್ಷರಗಳೆಂದು ಹೆಸರು.
\end{kbvrule}

\vspace{3pt plus 2pt}\noindent\hspace*{1.2em}\kbvlabel{ಉದಾ}\par\vspace{1pt}

\begin{kbvderiv}
\vspace{4pt plus 3pt minus 1pt}
\noindent{\renewcommand{\arraystretch}{1.24}\setlength{\tabcolsep}{5pt}
\begin{tabularx}{\textwidth}{X X}
ಬ್ + ಜ್ = (ಬ್ಜ) & ಸ್ + ತ್ + ರ್ = (ಸ್ತ್ರ) \\
ಕ್ + ತ್ = (ಕ್ತ) & ಷ್ + ಟ್ + ರ್ = (ಷ್ಟ್ರ) \\
\end{tabularx}}
\vspace{4pt plus 3pt minus 1pt}
\end{kbvderiv}

\kbvexercise{ಅಭ್ಯಾಸ ಪತ್ರಿಕೆ ೨}

\vspace{2.5pt plus 3pt minus 0.5pt}\noindent\hspace*{1.2em}\kbvnum{(1)}\hspace{0.55em}ಕನ್ನಡ ಅಕ್ಷರಮಾಲೆಯಲ್ಲಿರುವ ಅಕ್ಷರಗಳ ಸಂಖ್ಯೆಯೆಷ್ಟು?\par

\vspace{2.5pt plus 3pt minus 0.5pt}\noindent\hspace*{1.2em}\kbvnum{(2)}\hspace{0.55em}ಸ್ವರ ಎಂದರೇನು? ಸ್ವರಗಳಲ್ಲಿ ಎಷ್ಟು ವಿಧ?\par

\vspace{2.5pt plus 3pt minus 0.5pt}\noindent\hspace*{1.2em}\kbvnum{(3)}\hspace{0.55em}ಅನುನಾಸಿಕ ಎಂದರೇನು? ಉದಾಹರಣೆಗಳ ಸಹಿತ ತಿಳಿಸಿ.\par
